\section{Premise Validation (PV) Scorer and NLI Mechanics}
\label{sec:c1_pv}

\subsection{PV configuration and defaults}
\begin{table}[H]
\centering
\begin{tabular}{@{}p{0.25\linewidth}p{0.19\linewidth}p{0.50\linewidth}@{}}
\toprule
Parameter & Default & Effect \\
\midrule
\texttt{model\_name} & \texttt{microsoft/deberta-v3-base-mnli} & HuggingFace NLI checkpoint. \\
\texttt{device} & \texttt{auto} & Chooses CUDA if available, else CPU. \\
\texttt{batch\_size} & 8 & Number of premise-hypothesis pairs per forward pass. \\
\texttt{max\_length} & 256 & Tokenizer truncation cap. \\
\texttt{use\_fp16} & \texttt{True} & Enables autocast on CUDA for memory/perf efficiency, with FP32 fallback for known unstable cases. \\
\bottomrule
\end{tabular}
\caption{\texttt{PVConfig} defaults in \texttt{src/component1/pv.py}.}
\label{tab:c1_pv_defaults}
\end{table}

\subsection{Correct NLI pair ordering (critical for correctness)}
For each evidence item $e$ and claim $\Claim$:
\begin{itemize}
\item \textbf{premise} $=$ verbalized evidence text,
\item \textbf{hypothesis} $=$ claim text.
\end{itemize}
In tokenizer call order:
\begin{lstlisting}[style=jsonstyle]
tokenizer(evidence_texts, [claim_text] * len(evidence_texts), ...)
\end{lstlisting}
Swapping this order changes model behavior and invalidates score semantics.

\subsection{Label mapping strategy}
Different NLI checkpoints may index labels differently.
\texttt{PVScorer} determines indices in two stages:
\begin{enumerate}
\item inspect \texttt{model.config.id2label} and detect labels containing
\texttt{entail}, \texttt{neutral}, and \texttt{contra},
\item if incomplete, fallback to MNLI standard ordering
\([\texttt{contradiction},\texttt{neutral},\texttt{entailment}]\).
\end{enumerate}
Fallback map in code:
\begin{equation}
\texttt{entail}=2,\quad \texttt{neutral}=1,\quad \texttt{contra}=0.
\end{equation}

\subsection{Probability and derived signals}
For logits vector $\mathbf{z}(e)$:
\begin{equation}
\Pr(y=i \mid e,\Claim)=\frac{\exp(z_i)}{\sum_j \exp(z_j)}.
\end{equation}
Then:
\begin{align}
\Ent(e) &= \Pr(y=\text{entail}), \\
\Con(e) &= \Pr(y=\text{contra}), \\
\Neu(e) &= \Pr(y=\text{neutral}), \\
\Rel(e) &= \Ent(e)+\Con(e), \\
\Pol(e) &= \Ent(e)-\Con(e).
\end{align}

Ranges:
\begin{equation}
\Rel(e)\in[0,1],\quad \Pol(e)\in[-1,1].
\end{equation}
Interpretation:
\begin{itemize}
\item high $\Rel(e)$ means decision-relevant (non-neutral),
\item positive $\Pol(e)$ means support-leaning,
\item negative $\Pol(e)$ means contradiction-leaning.
\end{itemize}

\subsection{Batched inference procedure}
\begin{algorithm}[H]
\caption{PV scoring with cache-aware batched inference}
\begin{algorithmic}[1]
\Require claim text $\Claim$, evidence list $\Pool$, PV model/cache configs
\For{each evidence item $e \in \Pool$}
  \State verbalize premise text $p_e$
  \State compute \texttt{evidence\_hash} $=\texttt{sha256}(p_e)$
  \State lookup cache by key tuple (claim hash, evidence hash, model, max length, verbalizer id)
\EndFor
\State collect uncached premises into batch list
\If{uncached list non-empty}
  \State score in mini-batches with tokenizer and NLI model
  \State assign \texttt{PVResult} to each uncached evidence item
  \State write new scores into cache
\EndIf
\State return fully scored evidence pool
\end{algorithmic}
\end{algorithm}

\subsection{Operational safeguards}
\begin{itemize}
\item \texttt{torch.no\_grad()} and \texttt{model.eval()} avoid training-side effects.
\item FP16 autocast is used only when \texttt{device.type == "cuda"} and \texttt{use\_fp16=True}; known unstable model families (e.g., DeBERTa-v3) are forced to FP32.
\item If an FP16 overflow is detected at runtime, the batch is retried once in FP32.
\item Initialization failures are wrapped as \texttt{ValueError} with model name context.
\item Empty evidence list returns empty result list without model call.
\end{itemize}
